\chapter{LITERATURE REVIEW}

\section{Introduction}

Unmanned Aerial Vehicle (UAV) communication systems have attracted
significant research attention due to their wide range of applications
in surveillance, disaster monitoring, environmental analysis,
agriculture, and military operations. Reliable communication between
UAVs and ground stations plays a crucial role in ensuring successful
mission execution. However, UAV communication networks face multiple
challenges such as limited bandwidth, packet loss, dynamic topology,
and communication delays.

Several researchers have proposed various solutions to improve UAV
communication efficiency. These include routing optimization,
congestion control techniques, message scheduling algorithms,
priority-based transmission, and machine learning-based communication
strategies. This chapter presents a detailed review of the existing
research work related to UAV communication systems, MAVLink protocol,
congestion control, scheduling methods, and intelligent communication
mechanisms.


\section{Review of Existing Work}

UAV communication networks are generally implemented using Flying
Ad Hoc Networks (FANETs) or UAV Ad Hoc Networks (UAANETs). These
networks consist of multiple UAV nodes communicating wirelessly
without fixed infrastructure\cite{gupta2016surveyuav}
presented a comprehensive survey on UAV communication networks and
highlighted major challenges such as bandwidth limitation, routing
instability, and communication reliability\cite{bekmezci2013flying} introduced the concept of Flying
Ad Hoc Networks (FANETs) and discussed their unique features
including high mobility, dynamic topology, and intermittent
connectivity.

Mozaffari et al.~\cite{mozaffari2019tutorial} provided a detailed tutorial on UAV communication systems and emphasized challenges
such as interference management, delay control, and energy
efficiency. Nawaz et al.~\cite{nawaz2021uav} further identified
network congestion and packet loss as significant problems affecting
UAV communication performance.

The MAVLink protocol is widely used for communication between UAVs
and Ground Control Stations. Koubaa et al.~\cite{koubaa2019micro}
provided a comprehensive survey of MAVLink protocol architecture
and message formats. Their study showed that MAVLink uses static
message transmission rates, which may lead to inefficient bandwidth
utilization under heavy traffic conditions. Cordill et
al.~\cite{cordill2025comprehensive} analyzed security challenges
in UAV systems and reported that communication overload increases
system vulnerability.

Routing techniques also play a significant role in maintaining
reliable UAV communication. Perkins~\cite{perkins1999adhoc}
introduced the concept of Ad Hoc networking and routing methods
for dynamic communication systems. Clausen and
Jacquet~\cite{clausen2003olsr} proposed the Optimized Link State
Routing (OLSR) protocol to improve route reliability in mobile
networks. Fotouhi et al.~\cite{fotouhi2019survey} presented a
survey on UAV cellular communication systems and discussed
challenges related to interference and resource allocation.

Congestion control mechanisms are widely used to improve network
performance and reduce packet loss. Li et al.~\cite{li2020congestion}
presented a comprehensive survey on congestion control techniques
used in wireless networks. Anwekar and
Phulre~\cite{anwekar2023analysis} studied congestion issues in
FANETs and reported that uncontrolled traffic significantly
increases packet loss and communication delay. Verma et
al.~\cite{verma2024adaptive} proposed an adaptive congestion
control technique that dynamically adjusts transmission rates
based on network conditions.

Message scheduling and prioritization techniques have also been
widely studied. Li et al.~\cite{li2018prioritizing} proposed a
prioritization-based message scheduling method to ensure timely
delivery of critical messages in UAV networks. Halder et
al.~\cite{halder2022dynamic} introduced a distributed task
scheduling technique that improves coordination among UAV nodes.
Yang et al.~\cite{yang2025intelligent} developed an intelligent
queue scheduling mechanism that enhances communication efficiency
by managing packet transmission order based on priority levels.

Machine learning techniques are increasingly used in UAV
communication systems to improve performance and adaptability.
Breiman~\cite{breiman2001randomforest} introduced the Random
Forest algorithm, which is widely used for classification
problems due to its robustness and accuracy. Khan et
al.~\cite{khan2020mlnetwork} reviewed machine learning-based
congestion control techniques and highlighted their potential
in improving network efficiency. Gupta et
al.~\cite{gupta2021mlrouting} applied machine learning methods
for routing optimization and demonstrated improved communication
performance in UAV networks.

\subsection{Chosen Baseline Paper for Scheduling and Aggregation}

For the scheduling and aggregation component of this project, the
most relevant free-access baseline paper is Nguyen et al.~\cite{nguyen2024acs},
``Adaptive Clustering and Scheduling for UAV-Enabled Data Aggregation.''
This paper is important because it focuses directly on UAV scheduling,
clustering, and data aggregation, which are the same design themes used
in the combined reinforcement learning model developed in this project.

The paper proposes an adaptive scheduling framework for UAV-enabled data
collection and aggregation. However, it does not integrate MAVLink message
prioritization, congestion-aware transmission control, and joint RL-based
scheduling/aggregation in a closed-loop execution pipeline. That gap makes
it a strong but incomplete baseline for the present work.

\begin{table}[h]
\centering
\small
\resizebox{\textwidth}{!}{%
\begin{tabular}{|p{3.1cm}|p{5.6cm}|p{5.6cm}|}
\hline
\textbf{Aspect} & \textbf{Baseline Paper: Nguyen et al. (2024)} & \textbf{Our Proposed Combined Model} \\ \hline
Focus & Adaptive clustering, scheduling, and data aggregation & MAVLink prioritization, congestion control, RL scheduling, and aggregation \\ \hline
Decision Method & Rule-based / heuristic scheduling & Learned Q-table policy with closed-loop feedback \\ \hline
Message Awareness & General UAV data aggregation & Message semantic priority from Random Forest classification \\ \hline
Congestion Handling & Scheduling adapts to collection context & Explicit congestion-aware rate control and queue feedback \\ \hline
Deployment Style & Single optimization layer & Two-stage closed-loop architecture with policy and simulation services \\ \hline
\end{tabular}%
}
\caption{Comparison Between the Chosen Baseline Paper and the Proposed Combined Model}
\label{table:baseline_comparison}
\end{table}

This baseline is therefore used as the main literature anchor for the
scheduling and aggregation contribution, while the proposed model extends
it by adding message-level prioritization and reinforcement learning-driven
joint control.


\section{Research Gaps Identified}

Based on the literature review, several research gaps have been
identified in existing UAV communication systems.

Most existing UAV communication systems rely on static transmission
rates, which do not adapt to varying network conditions. This leads
to inefficient bandwidth utilization and increased packet loss
during congestion. Although routing techniques improve path
selection, they do not address message rate control at the
application level.

Many scheduling techniques prioritize messages based on predefined
rules but do not dynamically regulate transmission rates according
to network congestion levels. Similarly, several congestion control
methods treat all messages equally without considering the
importance of individual messages.

Machine learning techniques have been applied in routing and
network optimization; however, limited research integrates
machine learning-based message classification with congestion-aware
rate control mechanisms in MAVLink-based UAV communication systems.

In addition, existing baseline work such as Nguyen et al.~\cite{nguyen2024acs}
does not cover the combined closed-loop control strategy used in this project.
The present work addresses this gap by integrating semantic message
prioritization, adaptive rate control, and RL-based scheduling and aggregation
within one framework.

Therefore, there is a strong need to develop an intelligent
communication framework that combines machine learning-based
message prioritization with adaptive congestion-aware rate control
to improve communication reliability and Quality of Service (QoS)
in UAV networks.